\chapter{ماشین تورینگ}

تا اینجا سه مدل محاسباتی دیدیم:
\begin{itemize}
\item \lr{DFA}: قطعی، بدون حافظه موقت؛
\item \lr{NFA}: غیرقطعی، بدون حافظه موقت؛
\item \lr{PDA}: غیرقطعی، با یک پشته نامحدود.
\end{itemize}
در این فصل به مدلی می‌رسیم که \textbf{مدل یک کامپیوتر همه‌منظوره} است: \term{ماشین تورینگ}{Turing Machine} که آلن تورینگ در سال ۱۹۳۶ پیشنهاد کرد. ماشین تورینگ قطعی است و یک نوار نامحدود (از یک طرف) دارد که می‌تواند روی آن بنویسد و بخواند. هر کاری که یک کامپیوتر واقعی انجام دهد، ماشین تورینگ هم می‌تواند انجام دهد.

\section{ماشین تورینگ چیست؟}

ماشین تورینگ این اجزا را دارد:
\begin{itemize}
\item یک مجموعه متناهی از حالت‌ها با یک حالت شروع و دو حالت خاص: \textbf{پذیرش} و \textbf{رد}.
\item یک \textbf{نوار} که از یک طرف نامحدود است و به خانه‌ها تقسیم شده. در ابتدا رشته ورودی در خانه‌های اول نوار است و بقیه خانه‌ها \textbf{خالی} ($\blank$) هستند.
\item یک \textbf{هد} که روی نوار قرار دارد، نماد زیرش را می‌خواند، می‌تواند آن را عوض کند و یک خانه به چپ یا راست حرکت کند.
\end{itemize}
\begin{ltrfig}
\begin{tikzpicture}[font=\small]
  \node[draw=jzMain,thick,rounded corners,minimum width=2.2cm,minimum height=1cm,fill=jzMain!6] (c) at (0,1.5) {control};
  \foreach \s [count=\i] in {a,b,a,b,\sqcup,\sqcup} {\node[cell] (t\i) at (\i*0.65-1.6,0) {$\s$};}
  \node at (3.0,0) {$\cdots$};
  \draw[->,thick,jzMain] (c.south) -- (t1.north);
\end{tikzpicture}
\end{ltrfig}

\textbf{تفاوت‌ها با ماشین‌های متناهی:}
\begin{enumerate}
\item ماشین تورینگ هم می‌تواند روی نوار \textbf{بنویسد} و هم از آن بخواند.
\item هد می‌تواند هم به \textbf{چپ} و هم به \textbf{راست} حرکت کند.
\item نوار \textbf{نامحدود} است.
\item حالت‌های پذیرش و رد \textbf{فوراً} اثر می‌کنند: به محض ورود به آن‌ها ماشین متوقف می‌شود (لازم نیست ورودی تمام شده باشد).
\end{enumerate}

\textbf{نمایش گذارها:} برچسب $a\to b,R$ روی پیکانی از $q_1$ به $q_2$ یعنی: «اگر در $q_1$ هستی و هد نماد $a$ را می‌خواند، آن را با $b$ عوض کن، هد را یک خانه به راست ببر و به $q_2$ برو». برچسب $a\to R$ کوتاه‌شده $a\to a,R$ است (نماد عوض نمی‌شود) و $a\to L$ هم به همین شکل.

\section{مثال: سه دسته نماد برابر}

در فصل ۱۴ دیدیم $A=\{a^nb^nc^n\st n\ge0\}$ مستقل از متن نیست. حالا یک ماشین تورینگ برای آن می‌سازیم.

\begin{example}[ماشین تورینگ برای $a^nb^nc^n$]
فرض کنید ورودی آن‌قدر طولانی است که نمی‌توانید همه‌اش را به خاطر بسپارید، اما اجازه دارید روی ورودی جلو و عقب بروید و روی آن علامت بگذارید. ماشین در هر دور یک $a$، یک $b$ و یک $c$ را با هم جفت می‌کند. برای اینکه بداند کدام نمادها قبلاً بررسی شده‌اند، آن‌ها را با $x$، $y$ و $z$ جایگزین می‌کند. اگر همه نمادها جایگزین شدند، همه چیز جفت شده و ماشین می‌پذیرد؛ اگر جایی ناهماهنگی پیدا کرد، رد می‌کند.
\begin{automaton}[node distance=2.5cm]
  \node[state,initial] (q0) {$q_0$};
  \node[state,above=1.9cm of q0] (q1) {$q_1$};
  \node[state,right=2.9cm of q1] (q2) {$q_2$};
  \node[state,right=2.9cm of q0] (q3) {$q_3$};
  \node[state,below=1.7cm of q0] (q4) {$q_4$};
  \node[state,accepting,right=2.9cm of q4] (qa) {$q_{\mathrm{acc}}$};
  \path (q0) edge node {$a\to x,R$} (q1)
        (q1) edge[loop above] node[align=center] {$a\to R$\\$y\to R$} (q1)
        (q1) edge node {$b\to y,R$} (q2)
        (q2) edge[loop above] node[align=center] {$b\to R$\\$z\to R$} (q2)
        (q2) edge node {$c\to z,L$} (q3)
        (q3) edge[loop right] node[align=left] {$a\to L,\ b\to L$\\$y\to L,\ z\to L$} (q3)
        (q3) edge node[above] {$x\to R$} (q0)
        (q0) edge node[left] {$y\to R$} (q4)
        (q4) edge[loop below] node[align=center] {$y\to R$\\$z\to R$} (q4)
        (q4) edge node[below] {$\sqcup\to R$} (qa)
        (q0) edge node[pos=0.65] {$\sqcup\to R$} (qa);
\end{automaton}
گذارهایی که رسم نشده‌اند به حالت رد می‌روند. $q_1$ به دنبال اولین $b$ علامت‌نخورده می‌گردد، $q_2$ به دنبال اولین $c$ و $q_3$ هد را به اولین نماد علامت‌نخورده برمی‌گرداند. اگر در $q_0$ به جای $a$ یک $y$ ببینیم، همه $a$ها جفت شده‌اند و $q_4$ بررسی می‌کند که دیگر $b$ یا $c$ علامت‌نخورده‌ای نمانده باشد.
\end{example}

\section{تعریف صوری}

\begin{definition}[ماشین تورینگ]
یک ماشین تورینگ هفت‌تایی $M=(Q,\Sigma,\Gamma,\delta,q_0,\qacc,\qrej)$ است که در آن $Q$، $\Sigma$ و $\Gamma$ مجموعه‌های متناهی‌اند و
\begin{enumerate}
\item $Q$ مجموعه حالت‌ها،
\item $\Sigma$ الفبای ورودی است و نماد خالی را ندارد ($\blank\notin\Sigma$)،
\item $\Gamma$ الفبای نوار است، $\blank\in\Gamma$ و $\Sigma\subseteq\Gamma$،
\item $\delta:(Q-\{\qacc,\qrej\})\times\Gamma\to Q\times\Gamma\times\{L,R\}$ تابع گذار،
\item $q_0\in Q$ حالت شروع،
\item $\qacc\in Q$ تنها حالت پذیرش، و
\item $\qrej\in Q$ تنها حالت رد است و $\qacc\neq\qrej$.
\end{enumerate}
\end{definition}

اگر $\delta(q,s)=(r,t,D)$، یعنی: در حالت $q$ وقتی هد نماد $s$ را می‌خواند، به حالت $r$ برو، $s$ را با $t$ عوض کن و هد را در جهت $D\in\{L,R\}$ حرکت بده. چون دامنه $\delta$ همه حالت‌ها جز پذیرش و رد است، از هر حالت دیگر برای هر نماد نوار دقیقاً یک گذار وجود دارد و از حالت‌های پذیرش و رد هیچ گذاری نیست؛ ماشین با ورود به آن‌ها متوقف می‌شود.

\begin{note}
در عمل تقریباً هیچ‌وقت توصیف صوری کامل یک ماشین تورینگ را نمی‌نویسیم، چون خیلی بزرگ می‌شود. بیشتر توصیف‌های \textbf{سطح بالا} می‌دهیم. اما باید به خاطر داشت هر توصیف سطح بالا در واقع خلاصه‌ای از یک توصیف صوری است که با حوصله می‌توان آن را کامل نوشت.
\end{note}

\section{محاسبه ماشین تورینگ}

\begin{itemize}
\item ماشین ورودی $w=w_1\cdots w_n\in\Sigma^*$ را در $n$ خانه اول نوار می‌گیرد و بقیه نوار خالی است. هد روی چپ‌ترین خانه است.
\item چون $\blank\notin\Sigma$، اولین $\blank$ روی نوار انتهای ورودی را نشان می‌دهد.
\item محاسبه طبق تابع گذار پیش می‌رود. اگر ماشین بخواهد هد را از انتهای چپ نوار به چپ ببرد، هد سر جایش می‌ماند.
\item محاسبه تا ورود به حالت پذیرش یا رد ادامه می‌یابد و آنجا متوقف می‌شود. اگر هیچ‌کدام رخ ندهد، ماشین \textbf{برای همیشه} ادامه می‌دهد.
\end{itemize}

\begin{definition}[پیکربندی]
در طول محاسبه سه چیز تغییر می‌کند: حالت فعلی، محتوای نوار و مکان هد. به ترکیب این سه یک \term{پیکربندی}{Configuration} می‌گوییم و آن را به شکل $u\,q\,v$ می‌نویسیم: حالت فعلی $q$ است، محتوای نوار $uv$ (و بعد از آن فقط $\blank$) و هد روی اولین نماد $v$ قرار دارد.
\end{definition}

\begin{example}
در اجرای ماشین مثال قبل روی $aabbcc$، پیکربندی $xxy\,q_1\,bzc$ یعنی نوار $xxybzc$ است، ماشین در $q_1$ است و هد روی $b$ (چهارمین خانه) قرار دارد.
\end{example}

\begin{definition}[حرکت یک‌گامی]
پیکربندی $C_1$ پیکربندی $C_2$ را \textbf{نتیجه می‌دهد} اگر ماشین بتواند در یک گام از $C_1$ به $C_2$ برود. برای $a,b,c\in\Gamma$، $u,v\in\Gamma^*$ و $q_i,q_j\in Q$:
\begin{itemize}
\item \textbf{حرکت به چپ:} اگر $\delta(q_i,b)=(q_j,c,L)$، آنگاه $ua\,q_i\,bv$ نتیجه می‌دهد $u\,q_j\,acv$. در انتهای چپ نوار، $q_i\,bv$ نتیجه می‌دهد $q_j\,cv$.
\item \textbf{حرکت به راست:} اگر $\delta(q_i,b)=(q_j,c,R)$، آنگاه $ua\,q_i\,bv$ نتیجه می‌دهد $uac\,q_j\,v$. در انتهای چپ هم $q_i\,bv$ نتیجه می‌دهد $c\,q_j\,v$.
\item \textbf{انتهای راست:} پیکربندی $ua\,q_i$ همان $ua\,q_i\,\blank$ است؛ یعنی فرض می‌کنیم هد روی یک خانه خالی است و قوانین بالا را به کار می‌بریم.
\end{itemize}
\end{definition}

\begin{definition}[پذیرش]
\begin{itemize}
\item \textbf{پیکربندی شروع} روی ورودی $w$ برابر $q_0\,w$ است.
\item در \textbf{پیکربندی پذیرش} حالت $\qacc$ و در \textbf{پیکربندی رد} حالت $\qrej$ است. این دو \textbf{پیکربندی توقف} هستند و پیکربندی دیگری را نتیجه نمی‌دهند.
\end{itemize}
ماشین $M$ ورودی $w$ را \textbf{می‌پذیرد} اگر دنباله‌ای از پیکربندی‌ها $C_1,C_2,\dots,C_k$ وجود داشته باشد که (۱) $C_1$ پیکربندی شروع $M$ روی $w$ باشد، (۲) هر $C_i$ نتیجه دهد $C_{i+1}$ و (۳) $C_k$ یک پیکربندی پذیرش باشد.
\end{definition}

\begin{example}
اجرای ماشین $a^nb^nc^n$ روی $aabbcc$ (با $\vdash$ برای «نتیجه می‌دهد» و $\vdash^*$ برای چند گام):
\[
\begin{array}{l}
q_0aabbcc\vdash xq_1abbcc\vdash xaq_1bbcc\vdash xayq_2bcc\vdash xaybq_2cc\vdash xayq_3bzc\\
\vdash^* q_3xaybzc\vdash xq_0aybzc\vdash xxq_1ybzc\vdash xxyq_1bzc\vdash xxyyq_2zc\\
\vdash xxyyzq_2c\vdash xxyyq_3zz\vdash^* xq_3xyyzz\vdash xxq_0yyzz\vdash xxyq_4yzz\\
\vdash^* xxyyzzq_4\blank\vdash xxyyzz\blank\, q_{\mathrm{acc}}.
\end{array}
\]
\end{example}

\section{زبان‌های تشخیص‌پذیر و تصمیم‌پذیر}

\begin{definition}[تشخیص‌پذیر]
مجموعه رشته‌هایی که $M$ می‌پذیرد، \textbf{زبان} $M$ است و با $L(M)$ نشان داده می‌شود. زبانی را \term{تورینگ-تشخیص‌پذیر}{Turing-recognizable} (یا \term{بازگشتی شمارش‌پذیر}{Recursively Enumerable}) می‌نامیم اگر یک ماشین تورینگ آن را تشخیص دهد.
\end{definition}

اجرای یک ماشین تورینگ روی یک ورودی سه نتیجه ممکن دارد: \textbf{پذیرش}، \textbf{رد} یا \textbf{حلقه} (هرگز متوقف نمی‌شود). ماشین می‌تواند یک ورودی را با ورود به $\qrej$ یا با حلقه زدن «نپذیرد». مشکل این است که از بیرون نمی‌توان فهمید ماشینی که هنوز متوقف نشده، در حلقه افتاده یا فقط کارش طول می‌کشد.

\begin{definition}[تصمیم‌گیرنده و تصمیم‌پذیر]
ماشین تورینگی که روی \textbf{همه} ورودی‌ها متوقف شود یک \term{تصمیم‌گیرنده}{Decider} است. تصمیم‌گیرنده‌ای که زبانی را تشخیص دهد، آن زبان را \textbf{تصمیم می‌گیرد}. زبانی \term{تصمیم‌پذیر}{Decidable} (یا \textbf{بازگشتی}) است اگر یک ماشین تورینگ آن را تصمیم بگیرد.
\end{definition}

هر زبان تصمیم‌پذیر تورینگ-تشخیص‌پذیر است. در فصل ۱۹ می‌بینیم که زبانی وجود دارد که تشخیص‌پذیر هست اما تصمیم‌پذیر نیست.

\begin{note}
تصمیم‌پذیری \textbf{ویژگی زبان} است، نه ویژگی یک ماشین خاص. فرض کنید تصمیم‌گیرنده $M$ زبان $L$ را تصمیم بگیرد. می‌توانیم $M$ را طوری تغییر دهیم که به جای رد کردن، حلقه بزند: یک حالت جدید $q_{\mathrm{loop}}$ اضافه کنیم، همه گذارهایی را که به $\qrej$ می‌روند به $q_{\mathrm{loop}}$ ببریم و در $q_{\mathrm{loop}}$ با هر نمادی به راست حرکت کنیم. ماشین جدید دیگر تصمیم‌گیرنده نیست، اما زبانش همان $L$ است. $L$ همچنان تصمیم‌پذیر است، چون \emph{یک} تصمیم‌گیرنده (یعنی $M$) دارد.
\end{note}

\begin{theorem}
هر زبان منظم تصمیم‌پذیر است.
\end{theorem}
\begin{proof}
فرض کنید \lr{DFA} $M=(Q,\Sigma,\delta,q_0,F)$ زبان منظم $A$ را تشخیص دهد. ماشین تورینگ $M'=(Q',\Sigma,\Gamma,\delta',q_0,\qacc,\qrej)$ را با $Q'=Q\cup\{\qacc,\qrej\}$ و $\Gamma=\Sigma\cup\{\blank\}$ می‌سازیم:
\[\delta'(q,t)=\begin{cases}
(\delta(q,t),\ t,\ R) & t\in\Sigma\\
(\qacc,\ \blank,\ R) & t=\blank\ \text{\rl{و}}\ q\in F\\
(\qrej,\ \blank,\ R) & t=\blank\ \text{\rl{و}}\ q\notin F.
\end{cases}\]
$M'$ ورودی را از چپ به راست می‌خواند و حالتش را دقیقاً مثل $M$ عوض می‌کند. وقتی به اولین $\blank$ (انتهای ورودی) می‌رسد، بر اساس اینکه $M$ در حالت پذیرش است یا نه، می‌پذیرد یا رد می‌کند. پس $L(M')=L(M)=A$. همچنین $M'$ بعد از هر نماد به راست می‌رود و با رسیدن به $\blank$ متوقف می‌شود، پس روی هر ورودی متوقف می‌شود و تصمیم‌گیرنده است.
\end{proof}

\section{مثال: توان‌های دو}

\begin{example}
ماشین تورینگ $M_2$ زبان $A=\{0^{2^n}\st n\ge0\}$ (رشته‌هایی از $0$ که طولشان توانی از ۲ است) را تصمیم می‌گیرد.

\textbf{توصیف سطح بالا.} روی ورودی $w$:
\begin{enumerate}
\item نوار را از چپ به راست طی کن و یکی در میان $0$ها را خط بزن (تعداد $0$ها نصف می‌شود).
\item اگر در مرحله ۱ نوار فقط یک $0$ داشت، \textbf{بپذیر}.
\item اگر در مرحله ۱ نوار بیش از یک $0$ داشت و تعداد $0$ها فرد بود، \textbf{رد کن}؛ چون عدد فرد بزرگ‌تر از ۱ توانی از ۲ نیست. (ماشین با حالتش زوج یا فرد بودن تعداد $0$های دیده‌شده را به خاطر می‌سپارد.)
\item هد را به انتهای چپ نوار برگردان.
\item به مرحله ۱ برو.
\end{enumerate}

\textbf{توصیف صوری.} $M_2=(Q,\Sigma,\Gamma,\delta,q_1,\qacc,\qrej)$ با $Q=\{q_1,q_2,q_3,q_4,q_5,\qacc,\qrej\}$، $\Sigma=\{0\}$، $\Gamma=\{0,x,\blank\}$ و تابع گذار زیر:
\begin{automaton}[node distance=2.6cm]
  \node[state,initial] (q1) {$q_1$};
  \node[state,right=2.8cm of q1] (q2) {$q_2$};
  \node[state,right=2.8cm of q2] (q3) {$q_3$};
  \node[state,above right=1.4cm and 1.0cm of q2] (q5) {$q_5$};
  \node[state,below=1.6cm of q3] (q4) {$q_4$};
  \node[state,below=1.6cm of q1] (qr) {$q_{\mathrm{rej}}$};
  \node[state,accepting,below=1.6cm of q2] (qa) {$q_{\mathrm{acc}}$};
  \path (q1) edge node {$0\to\sqcup,R$} (q2)
        (q1) edge node[left,align=right] {$\sqcup\to R$\\$x\to R$} (qr)
        (q2) edge[loop above] node {$x\to R$} (q2)
        (q2) edge node {$0\to x,R$} (q3)
        (q2) edge node[left] {$\sqcup\to R$} (qa)
        (q3) edge[loop right] node {$x\to R$} (q3)
        (q3) edge[bend left=12] node[right] {$0\to R$} (q4)
        (q4) edge[bend left=12] node[left] {$0\to x,R$} (q3)
        (q4) edge[loop right] node {$x\to R$} (q4)
        (q3) edge node[above right] {$\sqcup\to L$} (q5)
        (q5) edge[loop above] node[align=center] {$0\to L$\\$x\to L$} (q5)
        (q5) edge node[above left] {$\sqcup\to R$} (q2)
        (q4) edge[bend left=25] node[below] {$\sqcup\to R$} (qr);
\end{automaton}
در $q_1$ اولین $0$ با $\blank$ جایگزین می‌شود تا انتهای چپ نوار علامت‌گذاری شود. $q_3$ و $q_4$ زوج و فرد بودن را دنبال می‌کنند و $q_5$ هد را به ابتدای نوار برمی‌گرداند. نمونه اجرا روی $0000$:
\[
\begin{array}{l}
q_10000\vdash\blank q_2000\vdash\blank xq_300\vdash\blank x0q_40\vdash\blank x0xq_3\blank\vdash\blank x0q_5x\blank\\
\vdash\blank xq_50x\blank\vdash\blank q_5x0x\blank\vdash q_5\blank x0x\blank\vdash\blank q_2x0x\blank\vdash\blank xq_20x\blank\\
\vdash\blank xxq_3x\blank\vdash\blank xxxq_3\blank\vdash\blank xxq_5x\blank\vdash\blank xq_5xx\blank\vdash\blank q_5xxx\blank\\
\vdash q_5\blank xxx\blank\vdash\blank q_2xxx\blank\vdash\blank xq_2xx\blank\vdash\blank xxq_2x\blank\vdash\blank xxxq_2\blank\vdash\blank xxx\blank q_{\mathrm{acc}}.
\end{array}
\]
\end{example}

\begin{summary}
\begin{itemize}
\item ماشین تورینگ: نوار نامحدود خواندنی/نوشتنی، هد دوطرفه، حالت‌های پذیرش و رد با اثر فوری.
\item $\delta:(Q-\{\qacc,\qrej\})\times\Gamma\to Q\times\Gamma\times\{L,R\}$. پیکربندی $uqv$ وضعیت کامل ماشین است.
\item سه نتیجه: پذیرش، رد، حلقه. تشخیص‌پذیر: یک ماشین می‌پذیردش. تصمیم‌پذیر: یک ماشین که همیشه متوقف می‌شود آن را می‌پذیرد.
\item هر زبان منظم تصمیم‌پذیر است و $\{a^nb^nc^n\}$ و $\{0^{2^n}\}$ هم تصمیم‌پذیرند.
\end{itemize}
\end{summary}
